\documentclass[11pt]{article}
\usepackage{amsmath, amssymb, geometry, amsthm, graphicx, enumerate}
\geometry{letterpaper, margin=1in}

\title{\textbf{Homework 1.1: Introduction to Differential Equations}}
\author{Tashfeen Omran}
\date{September 2026}

\begin{document}

\maketitle
\tableofcontents
\vspace{1cm}

% --- PART 1 ---
\section{Part 1: Comprehensive Introduction, Context, and Prerequisites}

\subsection{Core Concepts}
A \textbf{Differential Equation (DE)} is a mathematical equation that relates one or more unknown functions and their derivatives. In application, functions usually represent physical quantities, derivatives represent their rates of change, and the differential equation defines a relationship between the two. 

To properly classify and eventually solve differential equations, we must understand their anatomy:
\begin{itemize}
    \item \textbf{Independent Variable:} The variable with respect to which differentiation is taking place (often time $t$ or spatial coordinate $x$).
    \item \textbf{Dependent Variable:} The function whose derivative is being taken (often $y$ or $x$). Its value depends on the independent variable.
    \item \textbf{Ordinary Differential Equation (ODE):} An equation containing derivatives of one or more dependent variables with respect to a \textit{single} independent variable. Example: $\frac{dy}{dx} + y = 0$.
    \item \textbf{Partial Differential Equation (PDE):} An equation containing partial derivatives of one or more dependent variables with respect to \textit{two or more} independent variables. Example: $\frac{\partial u}{\partial t} = \frac{\partial^2 u}{\partial x^2}$.
    \item \textbf{Order:} The order of a differential equation is the order of the \textit{highest} derivative that appears in the equation. For example, if an equation contains $y''$ but no higher derivatives, it is a second-order equation.
    \item \textbf{Linearity:} An $n$-th order ODE is \textit{linear} if it can be written in the form:
    \[ a_n(x)y^{(n)} + a_{n-1}(x)y^{(n-1)} + \dots + a_1(x)y' + a_0(x)y = g(x) \]
    Two critical conditions must hold for an equation to be linear:
    \begin{enumerate}
        \item The dependent variable $y$ and all its derivatives $y', y'', \dots, y^{(n)}$ are of the first degree (i.e., the power of each term is exactly $1$).
        \item The coefficients $a_0(x), \dots, a_n(x)$ depend \textit{only} on the independent variable $x$. (This implies there are no terms where $y$ is multiplied by $y'$, no terms like $y^2$, and no transcendental functions of $y$ like $\sin(y)$ or $e^y$).
    \end{enumerate}
\end{itemize}

\subsection{Historical Context and Motivation}
The study of differential equations began with Isaac Newton and Gottfried Wilhelm Leibniz in the late 17th century. Newton developed calculus primarily to explain the motion of celestial bodies and physical mechanics. He stated his Second Law of Motion as $F = ma$, which is fundamentally a second-order differential equation: $F = m \frac{d^2x}{dt^2}$. Mathematicians needed to solve these equations to predict planetary orbits, determine the flight paths of projectiles, and understand heat transfer. Thus, differential equations were born out of the necessity to model a world in continuous motion.

\subsection{Key Formulas}
While this introductory chapter is more about classification than solving, the general form of an $n$-th order linear ODE is crucial:
\[ a_n(x)\frac{d^n y}{dx^n} + a_{n-1}(x)\frac{d^{n-1} y}{dx^{n-1}} + \dots + a_1(x)\frac{dy}{dx} + a_0(x)y = g(x) \]
For kinematics (constant acceleration), which appears in this assignment, the essential distance formula is:
\[ d(t) = v_0 t + \frac{1}{2}a t^2 \]

\subsection{Prerequisites}
To master this topic, you must be comfortable with:
\begin{itemize}
    \item \textbf{Calculus I \& II:} Differentiation rules (product, quotient, chain rule) and understanding Leibniz notation ($\frac{dy}{dx}$) versus Lagrange notation ($y'$).
    \item \textbf{Algebraic Manipulation:} Isolating variables, solving systems of equations, and simplifying rational expressions.
    \item \textbf{Kinematics:} Basic physics formulas for distance, velocity, and acceleration.
\end{itemize}

% --- PART 2 ---
\section{Part 2: Detailed Homework Solutions}

\subsection*{Problem 1}
\textbf{Equation:} $2\frac{d^2x}{dt^2} + 3\frac{dx}{dt} + 9x = 6\cos(4t)$
\vspace{0.2cm}

\noindent\textbf{Step-by-step Analysis:}
\begin{enumerate}
    \item \textbf{Variables:} The derivative is taken with respect to $t$, so the independent variable is $t$. The function being differentiated is $x$, so the dependent variable is $x$.
    \item \textbf{ODE vs PDE:} The equation uses ordinary derivatives ($d$) rather than partial derivatives ($\partial$), and there is only one independent variable ($t$). Thus, it is an \textbf{Ordinary Differential Equation (ODE)}.
    \item \textbf{Order:} The highest derivative present is the second derivative, $\frac{d^2x}{dt^2}$. Therefore, the \textbf{order is 2}.
    \item \textbf{Linearity:} Let us check the dependent variable $x$ and its derivatives. The terms containing $x$ are $2x''$, $3x'$, and $9x$. All of these are to the first power. They are not multiplied together, and they are not inside any nonlinear functions (the cosine function only contains the independent variable $t$). Thus, the equation is \textbf{linear}.
\end{enumerate}
\textbf{Final Answer:} Linear ordinary differential equation. Order: $2$. Independent variable: $t$. Dependent variable: $x$.

\subsection*{Problem 2}
\textbf{Equation:} $\frac{dy}{dx} = \frac{y(3-8x)}{x(5-10y)}$
\vspace{0.2cm}

\noindent\textbf{Step-by-step Analysis:}
\begin{enumerate}
    \item \textbf{Variables:} Derivative with respect to $x$ means $x$ is independent; $y$ is dependent.
    \item \textbf{ODE vs PDE:} Only ordinary derivatives are present with respect to a single variable $x$. It is an \textbf{ODE}.
    \item \textbf{Order:} The highest derivative is $\frac{dy}{dx}$ (first derivative). The \textbf{order is 1}.
    \item \textbf{Linearity:} To check for linearity, we can rearrange the equation. Multiply both sides by the denominator:
    \[ x(5 - 10y) \frac{dy}{dx} = y(3 - 8x) \]
    Notice the term $-10xy \frac{dy}{dx}$. The dependent variable $y$ is multiplied by its derivative $\frac{dy}{dx}$. This violates the condition for linearity. Thus, the equation is \textbf{nonlinear}.
\end{enumerate}
\textbf{Final Answer:} Nonlinear ordinary differential equation. Order: $1$. Independent variable: $x$. Dependent variable: $y$.

\subsection*{Problem 3}
\textbf{Equation:} $y\left[1 + \left(\frac{dy}{dx}\right)^2\right] = C$, where $C$ is a constant.
\vspace{0.2cm}

\noindent\textbf{Step-by-step Analysis:}
\begin{enumerate}
    \item \textbf{Variables:} $x$ is independent; $y$ is dependent.
    \item \textbf{ODE vs PDE:} Only ordinary derivatives are used. It is an \textbf{ODE}.
    \item \textbf{Order:} The only derivative present is $\frac{dy}{dx}$. The \textbf{order is 1}.
    \item \textbf{Linearity:} Expanding the equation gives $y + y\left(\frac{dy}{dx}\right)^2 = C$. There are two violations of linearity here: the derivative $\frac{dy}{dx}$ is squared, and the dependent variable $y$ is multiplied by a function of its derivative. Therefore, it is \textbf{nonlinear}.
\end{enumerate}
\textbf{Final Answer:} Nonlinear ordinary differential equation. Order: $1$. Independent variable: $x$. Dependent variable: $y$.

\subsection*{Problem 4}
\textbf{Equation:} $\sqrt{1-y} \frac{d^2y}{dx^2} + 2x \frac{dy}{dx} = 0$
\vspace{0.2cm}

\noindent\textbf{Step-by-step Analysis:}
\begin{enumerate}
    \item \textbf{Variables:} $x$ is independent; $y$ is dependent.
    \item \textbf{ODE vs PDE:} Only ordinary derivatives are used. It is an \textbf{ODE}.
    \item \textbf{Order:} The highest derivative is $\frac{d^2y}{dx^2}$. The \textbf{order is 2}.
    \item \textbf{Linearity:} The first term is $\sqrt{1-y} \cdot y''$. The dependent variable $y$ is inside a nonlinear function (the square root), and it is serving as a coefficient to the second derivative. This is a severe violation of linearity. It is \textbf{nonlinear}.
\end{enumerate}
\textbf{Final Answer:} Nonlinear ordinary differential equation. Order: $2$. Independent variable: $x$. Dependent variable: $y$.

\subsection*{Problem 5}
\textbf{Equation:} $8\frac{d^4y}{dx^4} = x(1-x)$
\vspace{0.2cm}

\noindent\textbf{Step-by-step Analysis:}
\begin{enumerate}
    \item \textbf{Variables:} $x$ is independent; $y$ is dependent.
    \item \textbf{ODE vs PDE:} Ordinary derivatives only. It is an \textbf{ODE}.
    \item \textbf{Order:} The highest derivative is the fourth derivative, $\frac{d^4y}{dx^4}$. The \textbf{order is 4}.
    \item \textbf{Linearity:} The dependent variable $y$ only appears in the term $8y^{(4)}$. The derivative is to the first power, and its coefficient ($8$) is a constant (which qualifies as a function of $x$). The right side of the equation only contains the independent variable $x$. Therefore, it is \textbf{linear}.
\end{enumerate}
\textbf{Final Answer:} Linear ordinary differential equation. Order: $4$. Independent variable: $x$. Dependent variable: $y$.

\subsection*{Problem 6}
\textbf{Prompt:} Write a differential equation that fits the physical description: "The rate of change of the population of a certain species at time $t$ is proportional to the first power of its population $M$."
\vspace{0.2cm}

\noindent\textbf{Step-by-step Analysis:}
\begin{enumerate}
    \item Identify the variables: Time is $t$ (independent), Population is $M$ (dependent).
    \item "Rate of change of the population at time $t$" translates to the derivative of $M$ with respect to $t$: $\frac{dM}{dt}$.
    \item "Is proportional to" means it equals some constant $k$ times the following expression.
    \item "The first power of its population $M$" translates simply to $M^1$ or $M$.
    \item Combining these parts: $\frac{dM}{dt} = kM$.
\end{enumerate}
\textbf{Final Answer:} The differential equation is $\frac{dM}{dt} = kM$.

\subsection*{Problem 7}
\textbf{Prompt:} Write a differential equation that fits the physical description: "The rate of change in the temperature $T$ of coffee at time $t$ is proportional to the difference between the temperature $M$ of the air at time $t$ and the temperature of the coffee at time $t$."
\vspace{0.2cm}

\noindent\textbf{Step-by-step Analysis:}
\begin{enumerate}
    \item Identify the variables: Time is $t$ (independent), Temperature of coffee is $T$ (dependent). The air temperature is $M$.
    \item "Rate of change in the temperature $T$" translates to $\frac{dT}{dt}$.
    \item "Is proportional to" means equals a constant $k$ times...
    \item "The difference between the temperature $M$ of the air... and the temperature of the coffee..." translates to $(M - T)$.
    \item Combining these: $\frac{dT}{dt} = k(M - T)$.
\end{enumerate}
\textbf{Final Answer:} The differential equation is $\frac{dT}{dt} = k(M - T)$.

\subsection*{Problem 8}
\textbf{Equation:} $\frac{\partial N}{\partial t} = \frac{\partial^2 N}{\partial r^2} + \frac{1}{r}\frac{\partial N}{\partial r} + kN$
\vspace{0.2cm}

\noindent\textbf{Step-by-step Analysis:}
\begin{enumerate}
    \item \textbf{Variables:} The derivative symbols are $\partial$, meaning partial derivatives. The derivatives are taken with respect to $t$ and $r$. Thus, $r$ and $t$ are the independent variables, and $N$ is the dependent variable.
    \item \textbf{ODE vs PDE:} Because there are multiple independent variables and partial derivatives are used, it is a \textbf{Partial Differential Equation (PDE)}.
    \item \textbf{Order:} The highest derivative is $\frac{\partial^2 N}{\partial r^2}$. The \textbf{order is 2}.
    \item \textbf{Linearity (Implicit in PDEs):} The dependent variable $N$ and all its partial derivatives appear to the first power. They are not multiplied together, nor are they inside nonlinear functions. Thus, it is a \textbf{linear} PDE.
\end{enumerate}
\textbf{Final Answer:} Partial differential equation. Order: $2$. Independent variables: $r$ and $t$. Dependent variable: $N$.

\subsection*{Problem 9}
\textbf{Prompt:} Two drivers, Alison and Kevin, are participating in a race. Beginning from a standing start, they each proceed with a constant acceleration. Alison covers the last $\frac{1}{7}$ of the distance in 3 seconds, whereas Kevin covers the last $\frac{1}{3}$ of the distance in 7 seconds. Who wins and by how much time?
\vspace{0.2cm}

\noindent\textbf{Step-by-step Solution:}
\begin{enumerate}
    \item \textbf{Set up the model:}
    Let $D$ be the total distance of the race.
    Because they start from rest, initial velocity $v_0 = 0$. The distance covered over time $t$ at constant acceleration $a$ is $d(t) = \frac{1}{2}at^2$.
    
    \item \textbf{Analyze Alison's race:}
    Let $t_A$ be Alison's total time, and $a_A$ be her acceleration.
    Total distance is $D = \frac{1}{2}a_A t_A^2$. \quad \textbf{(Eq. 1)}
    She covers the last $\frac{1}{7}$ of the distance in $3$ seconds. This means she covers the first $\frac{6}{7}$ of the distance in $(t_A - 3)$ seconds.
    So, $\frac{6}{7}D = \frac{1}{2}a_A (t_A - 3)^2$. \quad \textbf{(Eq. 2)}
    
    Divide Eq. 2 by Eq. 1 to eliminate $a_A$ and $D$:
    \[ \frac{\frac{6}{7}D}{D} = \frac{\frac{1}{2}a_A (t_A - 3)^2}{\frac{1}{2}a_A t_A^2} \implies \frac{6}{7} = \frac{(t_A - 3)^2}{t_A^2} \]
    Take the positive square root of both sides (since time must be positive and $t_A > 3$):
    \[ \sqrt{\frac{6}{7}} = \frac{t_A - 3}{t_A} \implies \sqrt{\frac{6}{7}} = 1 - \frac{3}{t_A} \]
    Isolate $t_A$:
    \[ \frac{3}{t_A} = 1 - \sqrt{\frac{6}{7}} \implies t_A = \frac{3}{1 - \sqrt{6/7}} \]
    Calculating the numerical value: $\sqrt{6/7} \approx 0.92582$, so $t_A \approx \frac{3}{1 - 0.92582} \approx 40.442$ seconds.

    \item \textbf{Analyze Kevin's race:}
    Let $t_K$ be Kevin's total time, and $a_K$ be his acceleration.
    Total distance is $D = \frac{1}{2}a_K t_K^2$. \quad \textbf{(Eq. 3)}
    He covers the last $\frac{1}{3}$ of the distance in $7$ seconds. This means he covers the first $\frac{2}{3}$ of the distance in $(t_K - 7)$ seconds.
    So, $\frac{2}{3}D = \frac{1}{2}a_K (t_K - 7)^2$. \quad \textbf{(Eq. 4)}
    
    Divide Eq. 4 by Eq. 3:
    \[ \frac{2}{3} = \frac{(t_K - 7)^2}{t_K^2} \implies \sqrt{\frac{2}{3}} = 1 - \frac{7}{t_K} \]
    Isolate $t_K$:
    \[ \frac{7}{t_K} = 1 - \sqrt{\frac{2}{3}} \implies t_K = \frac{7}{1 - \sqrt{2/3}} \]
    Calculating the numerical value: $\sqrt{2/3} \approx 0.81650$, so $t_K \approx \frac{7}{1 - 0.81650} \approx 38.146$ seconds.

    \item \textbf{Compare times:}
    Kevin's time is $\approx 38.146$ seconds. Alison's time is $\approx 40.442$ seconds.
    Because $38.146 < 40.442$, Kevin wins the race.
    The time difference is $t_A - t_K = 40.442 - 38.146 = 2.296$ seconds.
    Rounding to the nearest tenth, Kevin wins by $2.3$ seconds.
\end{enumerate}
\textbf{Final Answer:} Kevin wins the race by 2.3 seconds.

% --- PART 3 ---
\section{Part 3: In-Depth Analysis of Problems and Techniques}

\subsection{A) Problem Types and General Approach}
This homework assignment tests three distinct competencies:
\begin{enumerate}
    \item \textbf{Equation Anatomy and Classification (Problems 1-5, 8):} 
    \textit{Strategy:} Isolate the fundamental components. First, find all derivatives. The symbols ($d$ vs $\partial$) tell you if it is an ODE or PDE. The variables in the denominator of the derivative operator are independent; the variables in the numerator are dependent. Find the highest derivative for the order. Finally, scan the equation for any nonlinear occurrences of the \textit{dependent} variable.
    \item \textbf{Mathematical Translation / Modeling (Problems 6-7):}
    \textit{Strategy:} Treat the English sentence like a direct code. "Rate of change of $y$" translates directly to $y'$. "Proportional to" translates to $= k \cdot (\dots)$. 
    \item \textbf{Kinematic Root-Finding (Problem 9):}
    \textit{Strategy:} This is fundamentally an algebra problem masked in physics. The key strategy is setting up a system of equations where the unknown constants ($a$ and $D$) cancel out via division, leaving an equation with only one variable (time $t$).
\end{enumerate}

\subsection{B) Key Algebraic and Calculus Manipulations}
\begin{itemize}
    \item \textbf{Denominators as Nonlinear Triggers (Problem 2):} When a dependent variable is in the denominator, e.g., $\frac{1}{y}$, it is equivalent to $y^{-1}$. Since the power is not $1$, it is nonlinear. We proved this by cross-multiplying, which turned the term into $y \cdot y'$, clearly showing nonlinearity.
    \item \textbf{Division to Eliminate Variables (Problem 9):} In the kinematics problem, we had equations like $D = \frac{1}{2} a t^2$ and $\frac{6}{7}D = \frac{1}{2} a (t-3)^2$. We did not know $D$ or $a$. By dividing the second equation by the first, $\frac{D}{D}$ became $1$, and $\frac{a}{a}$ became $1$. This is a powerful algebraic trick for proportional physics problems.
    \item \textbf{Square Roots of Ratios (Problem 9):} When arriving at $\frac{(t-3)^2}{t^2} = \frac{6}{7}$, we simplified by taking the square root of both sides. This is allowed because we know time $t > 0$. It reduced a difficult quadratic equation into a simple linear algebraic equation.
\end{itemize}

% --- PART 4 ---
\section{Part 4: "Cheatsheet" and Tips for Success}

\subsection*{Summary of Key Rules}
\begin{itemize}
    \item \textbf{ODE vs PDE:} $d \implies$ ODE. $\partial \implies$ PDE.
    \item \textbf{Order:} Highest derivative power. (Do not confuse with the algebraic degree/power of the derivative!).
    \item \textbf{Linearity Tests for variable $y$:}
        \begin{itemize}
            \item Good: $y, y', y'', \cos(x)y, e^x y'$.
            \item Bad (Nonlinear): $y^2, y \cdot y', \frac{1}{y}, \sin(y), e^y, \sqrt{y}$.
        \end{itemize}
\end{itemize}

\subsection*{Cheats and Shortcuts}
\begin{itemize}
    \item \textbf{The $x$-immunity Rule:} A differential equation can be as incredibly complex as you want regarding the \textit{independent} variable (usually $x$ or $t$). Terms like $e^{x^2}\cos(\ln(x))$ do \textbf{not} make the equation nonlinear. Only the behavior of the dependent variable matters.
    \item \textbf{Proportionality Cheat:} The phrase "is proportional to" will always mean $= k \times (\dots)$. 
\end{itemize}

\subsection*{Common Pitfalls and Traps}
\begin{itemize}
    \item \textbf{Confusing Order and Degree:} $ \left(\frac{dy}{dx}\right)^4 = y$. The order is $1$ (first derivative), not $4$. The $4$ is the degree (algebraic power).
    \item \textbf{Hidden Non-linearity:} $\frac{dy}{dx} = \frac{x}{y}$. It looks linear at first glance because $y$ has no exponent, but multiplying by $y$ gives $y y' = x$, revealing it is nonlinear.
\end{itemize}

% --- PART 5 ---
\section{Part 5: Conceptual Synthesis and The "Big Picture"}

\subsection{A) Thematic Connections}
The core theme of differential equations is \textbf{Self-Referential Change}. In earlier calculus, you usually computed derivatives of explicitly known functions, like finding $f'(x)$ given $f(x) = x^2$. In this topic, the dynamic is reversed: you know a rule about how the function changes relative to its own current state, and you must work backward to discover the function itself. 

For instance, $y' = ky$ means "this function grows at a rate proportional to how big it currently is." You already know from Calculus I that the exponential function $e^{kx}$ is the only function that behaves this way. Thus, Differential Equations acts as the ultimate bridge between describing a physical law and discovering the mathematical function that obeys it.

\subsection{B) Forward and Backward Links}
\textbf{Backward:} This material builds entirely on the mechanical rules of differentiation and integration from Calculus I \& II. You must know your derivatives to recognize them here. 

\textbf{Forward:} Classifying equations is not just an academic exercise; it dictates the solution method. Just as you look at an integral to decide whether to use U-Substitution or Integration by Parts, determining if an ODE is "First-Order Linear" or "First-Order Separable" will tell you exactly which algorithm (e.g., Integrating Factors) you must use in Chapter 2. 

% --- PART 6 ---
\section{Part 6: Real-World Application and Modeling}

While differential equations model everything from springs to circuits, they are the absolute backbone of advanced modern finance.

\subsection{A) Concrete Scenarios in Finance and Economics}
\begin{enumerate}
    \item \textbf{Derivative Pricing (Black-Scholes):} In quantitative finance, determining the fair price of a European option requires understanding how its value changes as time passes and as the underlying stock price fluctuates. This is modeled by a famous Partial Differential Equation.
    \item \textbf{Interest Rate Modeling (The Vasicek Model):} Interest rates do not just grow to infinity; they tend to revert to a long-term average. Economists use a stochastic differential equation (a type of ODE that includes random noise) to model this "mean-reverting" behavior for bond pricing.
    \item \textbf{Portfolio Optimization (Merton's Problem):} An investor wants to maximize their expected lifetime wealth by continuously choosing how much to consume and how much to invest in a risky stock versus a risk-free bond. The solution requires solving the Hamilton-Jacobi-Bellman (HJB) equation, a highly complex nonlinear PDE.
\end{enumerate}

\subsection{B) Model Problem Setup: Black-Scholes PDE}
Let us model the pricing of a financial derivative (like a call option). We want to find the value of the option, $V$, which depends on two independent variables: the current price of the underlying stock ($S$) and time ($t$). 
\begin{itemize}
    \item \textbf{Dependent Variable:} $V(S, t)$ (The price of the option).
    \item \textbf{Independent Variables:} $S$ (Stock price) and $t$ (time).
    \item \textbf{Constants:} $r$ (risk-free interest rate) and $\sigma$ (volatility of the stock).
\end{itemize}
The fundamental principle of no-arbitrage in finance leads to the Black-Scholes PDE:
\[ \frac{\partial V}{\partial t} + \frac{1}{2}\sigma^2 S^2 \frac{\partial^2 V}{\partial S^2} + rS\frac{\partial V}{\partial S} - rV = 0 \]
Based on our homework, we can classify this! Because there are partial derivatives ($\partial$) with respect to $S$ and $t$, it is a \textbf{PDE}. The highest derivative is $\frac{\partial^2 V}{\partial S^2}$, so it is \textbf{Second-Order}. Finally, the dependent variable $V$ and all its derivatives appear to the first power, making it a \textbf{Linear} PDE.

% --- PART 7 ---
\section{Part 7: Common Variations and Untested Concepts}

The provided homework tests classification, but it omits a crucial conceptual step typically found in Chapter 1: \textbf{Verifying Solutions to Initial Value Problems (IVPs).}

\subsection*{Missing Concept: Verifying a Solution}
Before you learn how to solve differential equations, you are often asked to verify that a given function is indeed a solution. 

\textbf{Example Problem:} Verify that $y = C e^{-2x}$ is a solution to the ODE $y' + 2y = 0$.
\vspace{0.2cm}

\noindent\textbf{Solution Process:}
\begin{enumerate}
    \item Take the proposed solution: $y = C e^{-2x}$.
    \item Find the required derivatives. Here we need $y'$. Using the chain rule: $y' = -2C e^{-2x}$.
    \item Substitute $y$ and $y'$ back into the original differential equation:
    \[ (-2C e^{-2x}) + 2(C e^{-2x}) = 0 \]
    \[ 0 = 0 \]
    Because substitution yields a true statement ($0=0$), $y = C e^{-2x}$ is indeed a valid solution.
\end{enumerate}

% --- PART 8 ---
\section{Part 8: Advanced Diagnostic Testing: "Find the Flaw"}

Below are 5 problems. Each contains a subtle error in the provided solution. Your task as a student is to spot the error and provide the correction.

\subsection*{Flawed Problem 1: Classification}
\textbf{Problem:} Classify the equation $y'' + \sin(x) y = 0$ (ODE/PDE, Order, Linearity). \\
\textbf{Flawed Solution:} 
This equation contains ordinary derivatives, so it is an ODE. The highest derivative is $y''$, so it is 2nd order. Because the equation contains the transcendental function $\sin(x)$, it violates linearity. Therefore, it is a nonlinear ODE.

\textbf{Identify the Flaw:} The flaw is claiming that $\sin(x)$ makes the equation nonlinear. \\
\textbf{Correction:} Linearity only depends on the \textit{dependent} variable $y$. The function $\sin(x)$ acts as a coefficient depending only on the independent variable $x$, which is perfectly legal for a linear equation. The correct classification is \textbf{linear}.

\subsection*{Flawed Problem 2: Determining Order}
\textbf{Problem:} Find the order of the differential equation $\left(\frac{dy}{dx}\right)^3 + y = x$. \\
\textbf{Flawed Solution:} 
We examine the derivatives in the equation. The derivative term $\frac{dy}{dx}$ is raised to the 3rd power. Because the highest power of any derivative present is 3, the order of this differential equation is 3.

\textbf{Identify the Flaw:} The flaw is confusing algebraic degree (the exponent) with the order of the derivative. \\
\textbf{Correction:} The order is the highest derivative taken, which is just the first derivative $\frac{dy}{dx}$. The exponent $3$ makes it nonlinear, but the order is \textbf{1}.

\subsection*{Flawed Problem 3: Systems vs PDEs}
\textbf{Problem:} Classify the system $\frac{dx}{dt} + \frac{dy}{dt} = 0$ as an ODE or a PDE. \\
\textbf{Flawed Solution:} 
The equation contains two dependent variables, $x$ and $y$. Because there is more than one variable involved in the differentiation process, this must be a Partial Differential Equation (PDE).

\textbf{Identify the Flaw:} The flaw is looking at the number of dependent variables rather than independent variables. \\
\textbf{Correction:} An equation is a PDE only if there are multiple \textit{independent} variables (variables in the denominator). Here, both derivatives are taken with respect to a single independent variable, $t$. Therefore, this is an \textbf{ODE} (specifically, a linear system of ODEs).

\subsection*{Flawed Problem 4: Translating Proportionality}
\textbf{Problem:} Write a differential equation for: "The rate of change of price $P$ is inversely proportional to the price $P$." \\
\textbf{Flawed Solution:} 
Rate of change of $P$ translates to $\frac{dP}{dt}$. Inversely proportional means we subtract it from a constant. Thus, the equation is $\frac{dP}{dt} = k - P$.

\textbf{Identify the Flaw:} The flaw is a fundamental misunderstanding of the word "inversely proportional." \\
\textbf{Correction:} Inversely proportional means proportional to the reciprocal (division), not subtraction. The correct equation is $\frac{dP}{dt} = \frac{k}{P}$.

\subsection*{Flawed Problem 5: Identifying Nonlinear Forms}
\textbf{Problem:} Classify $x^2 y'' + y y' = 0$ for linearity. \\
\textbf{Flawed Solution:} 
We look at the dependent variable $y$. The first term has $y''$, which is to the first power. The second term is $y y'$, where $y$ is to the first power and $y'$ is to the first power. Since there are no powers greater than 1, and no functions like $\cos(y)$, the equation is linear.

\textbf{Identify the Flaw:} The flaw is ignoring the product rule of linearity; the variables cannot be multiplied by one another. \\
\textbf{Correction:} The term $y y'$ is a product of the dependent variable and its derivative. A linear ODE requires that coefficients of $y$ and its derivatives only depend on the independent variable $x$. Thus, the equation is \textbf{nonlinear}.

\end{document}